\documentclass[../main]{subfiles}
\begin{document}
\ifSubfilesClassLoaded{\setcounter{page}{27}}{}
\section{Загальна (суспільна) власність}
\label{sec:05_obshhaya_sobstvennost}

Перешкодою на шляху до комунізму є приватна власність. Перша праця, в якій Маркс серйозно та ґрунтовно розмірковує про комунізм, – це «Економічно-філософські рукописи 1844 року». Розділ про комунізм знаходиться у другому рукописі одразу після розділу про приватну власність. Він починається роздумами про приватну власність, про працю як сутність приватної власності та про розуміння приватної власності та праці у різних представників соціалізму та комунізму як ідейного напряму. Комуністи виступають проти приватної власності (в тому числі й до Маркса) на засоби виробництва, оскільки вона перетворює людей на рабів інших людей і породжує всі жахи капіталістичного суспільства.

Приватна та особиста власність – це не одне й те саме. Речі, що використовуються безпосередньо та призначені для особистого споживання, які належать їхньому власнику, за формою є приватною власністю, так само як і заводи, фабрики, магазини, корпорації, великі земельні ділянки. Цей факт затьмарює суть справи. Тому, коли антикомуністи чують про знищення приватної власності, найпростіші з них думають, що йдеться про об’єднання в загальну власність абсолютно всіх речей, і починають боятися за свої особисті речі. У цьому питанні ми можемо їх заспокоїти. Речі особистого вжитку не становлять інтересу для комуністів у тому сенсі, щоб у когось їх відібрати, а лише в тому, щоб усі були ними забезпечені в достатній мірі для повноцінного та різнобічно насиченого життя. Економічно ці речі перебувають не в приватній, а в особистій власності.

Знищення приватної власності не скасовує особистої власності (приватна власність – це відносини між людьми щодо речей, а не самі речі), оскільки споживання за необхідності має бути особистим. Навпаки, знищення приватної власності забезпечує доступ усіх членів суспільства до предметів, необхідних для особистого споживання, а там, де це більш раціонально, переводить споживання в суспільні форми. Наприклад, організація споживчих кооперативів дозволяє підвищити ефективність використання речей і зменшує потребу в їх виробництві та зберіганні, а отже, також і витрати суспільного часу на це. Так, організація простого споживчого кооперативу з використання предметів, які використовуються не так часто, дозволяє не захаращувати речами простір і максимально ефективно використовувати кожен із них аж до їх морального та фізичного зносу. Таким чином, річ переходить з рук в руки, і кожен учасник кооперативу може отримати від неї користь. Приклади споживчих кооперативів багато разів описувалися в літературі. Це легко реалізувати з найрізноманітнішими предметами, які перестають належати комусь особисто, але при цьому використовуються всіма в міру, необхідну для задоволення потреб. Комунізм в цьому випадку протистоїть приватній власності зовсім не з того боку, якого побоюються його противники. Суспільні форми споживання спрямовані не на те, щоб у когось щось відібрати, а на те, щоб надати доступ до предметів споживання всім.

Що стосується таких соціальних форм організації побуту, як спільне використання пральних машин, спільне приготування та доставка їжі, спільне використання прибиральної техніки, то тут економія ресурсів і людського часу (як безпосередньо, так і опосередковано – як економія часу на виробництво ресурсів) при тому ж рівні споживання може бути дуже значною. У кожний момент часу вона пропорційна можливостям техніки. Але навіть зараз, наприклад, приготування каші для однієї дитини в домашніх умовах і приготування каші для 200 дітей в умовах дитячого садка або молочної кухні, оснащених сучасною технікою, займає однаковий час.

Те саме стосується і приготування їжі для дорослих. Організоване харчування дозволяє використовувати більш ефективні технології, які неможливо застосувати в індивідуальних домашніх умовах, і забезпечує набагато вищий рівень автоматизації всіх процесів. При цьому можна краще контролювати якість, що зрештою вигідно всім учасникам процесу. Рух у цьому напрямку існує і за капіталізму. Сам капіталізм у свій спосіб створює та стандартизує систему громадського харчування (наприклад, McDonald's пропонує однакові картопляні чіпси в усіх куточках світу), пральні тощо. Але за капіталізму вони все одно не використовуються в громадський спосіб. Форми спільного споживання постійно стикаються з протидією в товарному виробництві, а виробництво, як відомо, визначає споживання.

Усі тягари побуту, в яких задіяні численні люди, можна замінити передовим виробництвом, що базується на сучасних технологіях; приватний транспорт, який заважає перехожим і створює затори, замінюється на добре розвинений громадський транспорт, і так далі у всіх сферах людського життя, де це пішло б на користь усім. Наприклад, у сучасних великих містах, де розвинений громадський транспорт, ті, хто ним користується, особливо в години пік, добираються до потрібного місця більш ефективно. У Москві діти, які їздять до школи на метро, здебільшого приїжджають вчасно, ті ж, кого підвозять на особистому автомобілі, як правило, спізнюються, оскільки спроби використання обмеженого міського простору в приватних цілях виявляються менш ефективними, ніж суспільне використання. Це також пояснює, чому міська влада під час снігопадів закликає громадян відмовитися від використання особистого транспорту та пересісти на громадський, щоб мінімізувати кількість аварійних ситуацій на дорогах.

Товарне виробництво сприяє розвитку індивідуальних форм споживання. Навіть попри їхню нераціональність та неефективність, вони є умовою постійного попиту на товари на ринку, умовою руху капіталу. Саме в цьому сенсі приватна власність перетинається з особистою і впливає на неї.

Приватна власність, або, точніше кажучи, приватна власність на засоби виробництва – це, перш за все, можливість отримання додаткового прибутку (додаткової вартості у формі прибутку або ренти), тобто експлуатація інших людей.

Комуністи говорять про об’єднання приватної, а не особистої власності – засобів суспільного виробництва, які зараз використовуються не на благо всього суспільства, а слугують експлуатації людини людиною. Приватна власність – це не речі, а суспільні відносини, що проявляються через речі.

Її суть становить експлуатація. Представники домарксистського комунізму розуміли знищення приватної власності як її поширення на все суспільство, перетворення всіх членів суспільства на робітників та зрівняння всіх як щодо праці, так і щодо споживання. Приватна власність розглядається багатьма представниками домарксистського комунізму лише як матеріальне багатство, яким можна володіти разом. Наприклад, Прудон розглядав приватну власність лише як капітал, який підлягає знищенню. Ближче до Маркса були ті теоретики, які вважали джерелом приватної власності та її згубного впливу на людину певний характер праці, «працю, яка зрівнялася, роздробилася і тому є невільною», однак вони зводили таку працю до певних видів праці. «Фурьє, подібно до фізіократів, знову ж таки вважає землеробську працю принаймні найкращим видом праці, тоді як, навпаки, Сен-Сімон вважає, що суть полягає в промисловій праці як такій, і тому прагне до беззастережного панування промисловців і поліпшення становища робітників. І, нарешті, комунізм є позитивним вираженням скасування приватної власності; на перших етапах він виступає як загальна приватна власність». «Позитивне» в даному випадку означає «збереження всіх багатств культури та діяльних здібностей людини, розвинених у суспільстві товарного виробництва».

Маркс піддає безжальній критиці розуміння заперечення приватної власності як її поширення на все суспільство. Він показує, що такий комунізм не долає приватну власність, а є лише теоретичним завершенням приватної власності як суспільного відносини. Він зазначає, що панування власності, як сукупності речей, над таким комунізмом, як ідейно-теоретичною позицією та суспільним рухом, залишається дуже сильним. Суспільні відносини все ще мисляться лише в категоріях приватної власності. Такий комунізм «прагне знищити все те, чим, на засадах приватної власності, не можуть володіти всі; він хоче насильницьки абстрагуватися від таланту тощо». Безпосереднє фізичне володіння постає перед ним як єдина мета життя та існування; категорія робітника не скасовується, а поширюється на всіх людей; ставлення до приватної власності залишається ставленням усього суспільства до світу речей; нарешті, цей рух, який прагне протиставити приватній власності загальну приватну власність, виражається в цілком звірячій формі, коли він протиставляє шлюбу (який, по суті, є певною формою ексклюзивної приватної власності) спільне користування жінками, де, отже, жінка стає суспільною та загальною власністю.

Для Маркса ідея спільного володіння жінками – це ставлення до людини як до речі, що відображає й ставлення людини до світу речей, створених людиною для людини. «Так само, як жінка тут переходить від шлюбу до загальної проституції, так і весь світ багатства, тобто матеріальна сутність людини, переходить від виключного шлюбу з приватним власником до універсальної проституції з усім суспільством». Маркс називає такий комунізм грубим. Під багатством тут розуміється сукупність усіх речей, створених людьми для людей, в яких втілені як діяльні здібності людського суспільства, так і потреби, які воно сформувало.

Теорія грубого комунізму, яка заперечує індивідуальність людини, лише абсолютизує в уявленні реальну тенденцію суспільства до приватної власності. Там, де багатство людських відносин зводиться до відносин приватної власності, індивідуальність зводиться до економічного агента приватної власності. Власне, особистість цих агентів (як пролетарія, так і капіталіста) абсолютно не має значення для руху приватної власності. Пролетарій цікавить капіталіста не як особистість, а виключно як продавець робочої сили певної якості, капіталіст – пролетарія – виключно як її покупець. Те саме відбувається в актах купівлі-продажу будь-яких інших товарів. Усі інші особистісні якості економічних агентів не мають відношення до руху приватної власності. Якщо люди протягом усього свого життя не виходять за межі цих відносин, їхня особистість починає зводитися лише до цих визначень.

Примітивний комунізм передбачає таке скасування приватної власності, яке не є справжнім її подоланням. Це видно з абстрактного заперечення досягнень культури, розгляду їх як надлишкових лише тому, що вони не потрібні пролетаріату як пролетаріату, і (за виразом Маркса) прагнення до неприродної простоти бідного, грубого і позбавленого потреб людства, яке не може подолати приватну власність, оскільки ще не доросло до неї.

Лозунг «Все відібрати і поділити», характерний для сучасних представників псевдокомунізму, стоїть ще нижче. Вони навіть не дійшли до об’єднання, хоча б у тому вигляді, в якому його пропонували представники примітивного комунізму – об’єднання в умовах злиднів.

Маркс протиставляє грубим комуністичним теоріям комунізм як позитивне скасування приватної власності. Такий комунізм, як історичний процес, зберігає багатство всієї попередньої історичної еволюції у вигляді діяльних здібностей, потреб, матеріальної та духовної культури людства. Комунізм робить усе це надбанням кожного члена суспільства. Як історичний процес він має свідомий, а не стихійний характер, оскільки суспільне виробництво передбачає свідоме ставлення членів суспільства до виробництва та відтворення свого суспільного життя. Комунізм передбачає освоєння об’єктивних законів суспільного розвитку, які людство може використати на свою користь так само, як закони фізики, хімії та біології в сучасній промисловості.

Уся попередня історія людства створює умови для такого комунізму, який долає приватну власність, але при цьому зберігає все багатство попереднього розвитку. У процесі всієї попередньої історії формується і відповідна свідомість – теорія комунізму. Комунізм розуміється як процес знищення приватної власності: «Для знищення ідеї приватної власності цілком достатньо ідеї комунізму. Для знищення ж приватної власності в реальній дійсності потрібна реальна комуністична дія». Варто зазначити один важливий, на нашу думку, момент.

Йдеться про вживання термінів «соціалізм» і «комунізм», усталене значення яких не сформувалося одразу. У наведених вище рукописах 1844 року соціалізмом ще називають суспільство, де приватна власність не скасовується, а вже скасована. Тоді як комунізм розуміється як процес (акт) заперечення (скасування) приватної власності. Саме тому тут він не розглядається як форма суспільства, а лише як енергійний принцип найближчого майбутнього.

Цей факт дозволяє зрозуміти, чому саме ці питання були підняті і саме так викладені Марксом і Енгельсом у «Маніфесті комуністичної партії», всього через 3 роки після написання «Рукописів». Більш чітко простежуються корені основного положення, викладеного в «Маніфесті»: «Сучасна буржуазна приватна власність є останнім і найповнішим вираженням такого виробництва і присвоєння продуктів, яке тримається на класових антагонізмах, на експлуатації одних іншими. У цьому сенсі комуністи можуть висловити свою теорію одним положенням: знищення приватної власності».

Це скасування, як вже зазначалося, передбачає збереження всього культурного багатства, активних здібностей і потреб, розвинених на основі приватної власності. Воно здійснюється не заради себе, а для існування комуністичного суспільства з відповідною власністю та способом спілкування.

Якою б не була комуністична спільнота, людина як природна істота повинна жити в гармонії з природою та привласнювати її: все, чим користується людина і що споживає, створено з природних матеріалів, за допомогою природних процесів, поставлених на службу людині, або є самими цими природними процесами. Через 20 років після написання «Маніфесту» у Вступі до «Економічних рукописів 1857–1859 років» Маркс писав: «Будь-яке виробництво є привласненням індивідуумом предметів природи в межах певної суспільної форми та за її допомогою. У цьому сенсі буде тавтологією сказати, що власність (привласнення) є умовою виробництва. Смішно, однак, робити звідси висновок про певну форму власності, наприклад, про приватну власність (що, до того ж, передбачає як умову протилежну форму — відсутність власності)».

Комуністичному суспільству відповідає загальна, тобто суспільна форма власності на всі суспільні засоби виробництва і, відповідно, загальне суспільне управління всім суспільним виробництвом як єдиним цілим. Суспільній природі виробництва, розвиненій на основі приватної власності, має відповідати суспільний характер присвоєння продуктів цього виробництва. Маркс і Енгельс писали про це десятки разів, у різних своїх працях. Мабуть, це єдина теза, яку запам’ятали сучасні комуністи, навіть якщо вони не розуміють її значення.

Для початку цього історичного руху в найбільш розвинених країнах Маркс і Енгельс запропонували такі заходи: «1. Експропріація земельної власності та спрямування земельної ренти на покриття державних витрат. 2.

Високий прогресивний податок. 3. Скасування права на спадок. 4. Конфіскація майна всіх емігрантів і бунтівників. 5. Централізація кредиту в руках держави шляхом створення національного банку з державним капіталом і з винятковою монополією. 6. Централізація всього транспорту в руках держави. 7. Збільшення кількості державних фабрик, засобів виробництва, розчищення земель під обробіток і поліпшення ґрунтів за загальним планом. 8. Обов’язкова праця для всіх, створення промислових армій, особливо для сільського господарства. 9. Поєднання сільського господарства з промисловістю, сприяння поступовому усуненню відмінностей між містом і селом. 10. Загальна і безкоштовна освіта для всіх дітей. Усунення фабричної праці дітей у її сучасному вигляді.

Як бачимо, з точки зору поставленої задачі щодо ліквідації приватної власності, ці заходи не здаються надто радикальними, але загалом вони відповідають вимогам маніфесту. Маніфести мають бути чіткими, зрозумілими та визначати основний напрямок розвитку. Тим паче, самі автори стверджували, що заходи залежать від конкретно-історичних умов, а тому можуть бути й іншими. Про заходи ми спеціально поговоримо в одному з наступних випусків, розглядаючи їх через призму досвіду соціалізму XX століття.

А зараз слід сказати ще кілька слів про те, що таке приватна власність, і що саме комуністи ставлять собі за мету ліквідувати, зберігаючи при цьому все позитивне, щоб «замість старого буржуазного суспільства з його класами та класовими протиріччями» постала «асоціація, в якій вільний розвиток кожного є умовою вільного розвитку всіх».

Приватна власність – це продукт невільної праці і водночас засіб виробництва невільного характеру праці: такої праці, коли ні процес, ні продукт праці безпосередньо не належать колективу працівників, праці, яка сама по собі позбавлена сенсу для працівників, а є лише засобом для досягнення зовнішніх цілей (грошей). Сутність приватної власності – у такому характері праці. Вона досягає меж свого розвитку в капіталістичному суспільстві розвиненого товарного виробництва, завдяки доведенню до крайньої міри суспільного поділу праці, який перетворює продукти праці на товари (у товарному виробництві).

Суспільний поділ праці протягом тривалого часу сприяв розвитку людства. Він підвищував продуктивність праці та створював матеріальну базу для того, щоб передати машиноподібні функції, які змушений виконувати працівник через недосконалість інструментів, самим цим інструментам – машинам. Але він не лише об’єднує людей, мовляв, Вася – слюсар, а Петя – програміст, кожен виконує свою справу, а над ними ще стоїть фахівець-менеджер, який керує ними для спільного блага. Такої мети не ставить перед собою жоден із учасників процесу. Усі люди в такій ситуації опиняються пов’язаними між собою зовнішнім чином – через продукти своєї праці. Вони змушені обмінювати свій продукт на той, якого у них немає, залишаючись у процесі праці байдужими один до одного. Неважливо, хто і як робить цей продукт, важливо, що цей продукт чужої праці потрібен мені.

Суспільний поділ праці пройшов декілька важливих етапів, які відіграли позитивну роль у розвитку людства. Про них варто розповісти в окремій праці, присвяченій саме цьому питанню. Тут важливо лише зазначити, що суспільний розвиток – це продукт історії. Марксисти відкидають усі забобони щодо біологічної, природної пристосованості членів суспільства до певних видів праці. Ці забобони не витримують навіть поверхневої наукової критики. Вони не відповідають не лише тому, що відомо про історію цього питання, але й тій історії, яка відбувається на наших очах. Щороку з’являються та зникають численні нові професії. Батьки часто не розуміють, у чому суть роботи їхніх дітей, і, звісно, не могли передати до неї схильність у спадок разом із генами.

Відколи праця стає частковою, а продукти людської праці перетворюються на товари, тобто починають виготовлятися спеціально для обміну, ці продукти починають домінувати над людьми. Речі – продукти їхньої власної праці – і гроші, як представники всіх речей, набувають влади над живими людьми. А відколи робоча сила стає товаром, таке домінування набувають люди, у яких є достатньо грошей, щоб купувати засоби виробництва та робочу силу. Ці люди володіють живою людською працею. Економічно вони є лише представниками капіталу. Їхня функція – приєднувати до капіталу чужу працю, щоб збільшилася вартість капіталу. Зрештою, це призводить до того, що майже всі члени суспільства так чи інакше працюють заради збільшення капіталу – самозростання вартості. «Лише продукти самостійної, незалежної одна від одної приватної праці протистоять один одному як товари». Тому знищення приватної власності – це знищення товарного виробництва взагалі, а не лише капіталістичного товарного виробництва. Вартість прагне до самозбільшення.

Це закон, який регулює суспільство, де кожен обмежений своєю частковою діяльністю, або, іншими словами, суспільство розвиненого товарного виробництва. Закони сучасних держав – це природне юридичне вираження цього об’єктивного закону виробництва, заснованого на невільній і роздробленій праці.

До певного моменту прагнення капіталу до самозбільшення є рушійною силою розвитку суспільства. Але воно породжує суспільні сили, яким стає тісно в рамках капіталу. Використання цих сил виключно в цілях збільшення капіталу призводить до катастрофічних наслідків.

Ці сили потрібно спрямовувати, і це можливо зробити на благо суспільства лише разом, для розвитку всіх, за загальним розумним планом. А для цього необхідно розвивати універсальні людські здібності (опанувати універсальну чутливість, мислення та моральність), виходячи за межі своєї професії, що не дозволяє розібратися навіть у «сусідній» професії, а тим більше – у суспільних процесах загалом. Тому знищення приватної власності дорівнює знищенню невільної форми праці. Це відбувається в міру розвитку вільної комуністичної діяльності та перетворення її на надбання всіх людей.

Комунізм протиставляє приватній власності загальну власність як форму спільної практичної взаємодії людей із засобами виробництва та природними ресурсами. Загальна (суспільна) власність забезпечує найбільш ефективне використання виробничих можливостей не лише для суспільства в цілому, але й для кожного його члена. Вона передбачає активну участь і активну волю щодо суспільних ресурсів від кожного члена суспільства. Відповідно, якщо приватна власність є продуктом поділу праці (розділення її на все більш приватні види праці)\footnote{Саме тому, що в СРСР суспільний поділ праці не було подолано (тенденція до його подолання не стала домінуючою), а пролетаріат відтворювався як пролетаріат, тобто як клас буржуазного суспільства, відбулася повна реставрація панування приватної власності.} і є цим поділом, то загальна власність передбачає, як уже зазначалося вище, об’єднання праці, усунення її поділу. Суспільна власність передбачає вільну людську діяльність за критеріями істини, добра та краси. Останнє є справжньою економічною основою для реального, а не лише формального, об’єднання та справжньої збіжності приватних і суспільних інтересів усіх членів суспільства.

Під загальною або суспільною власністю не слід розуміти власність колективних господарств і кооперативів, або інших груп людей, які виробляють товари та торгують ними з усім іншим світом. Вони не виходять за межі приватної власності, оскільки не виходять за межі товарного виробництва. З суспільством члени таких господарств пов’язані як власники товарів і виробники товарів. Сам характер товарного виробництва визначає й відносини в цих колективах, оскільки ставить виробництво вартості в основу, робить прибуток метою роботи колективу.

Зачатки спільної власності також не варто шукати, наприклад, у власності сім’ї. У XIX столітті деякі соціалісти уявляли собі, що суспільна власність може виникнути шляхом «розширення» сімейної власності до масштабів усього суспільства. Суспільство уявлялося у вигляді великої сім’ї з притаманними їй способами розподілу. Але сім’я – це втілення приватного характеру присвоєння продуктів праці. Вона сама, у своїй сучасній формі, є продуктом розвитку приватної власності. Настільки, наскільки в сім’ї здійснюється виробництво та відтворення людини як робочої сили (людина як економічний агент товарного суспільства), вона цілком вписується і повністю відтворюється за законами розвитку приватної власності. У тій мірі, в якій вона відтворює всіх своїх членів і є основною споживчою «клітиною», вона цілком і повністю визначається капіталістичним виробництвом. Це економічно передбачається як у формі заробітної плати, так і у формі прибутку. Більше того, сучасна сім’я як господарська одиниця не лише не виключає, але й прямо передбачає експлуатацію людини людиною, що несумісне з комунізмом. Тому власність сім’ї не можна вважати основою суспільної власності – власності, характерної для комунізму.

І тим більше, комуністичну власність, яка могла б стати ключовою точкою на шляху до комунізму, не можна побачити в жодних ізольованих спільнотах, не пов’язаних із суспільним виробництвом (наприклад, у релігійних громадах, що базуються на натуральному господарстві).

Але де ж тоді можна побачити суспільну власність, характерну для комунізму? Очевидно, що лише там, де існує вільна людська діяльність у її безпосередньо-колективній формі. Така діяльність, розглядаючи її з боку використання та розпорядження ресурсами, з боку привласнення за допомогою суспільної форми, і є дійсним втіленням суспільної власності.

У попередній главі ми навели лише кілька прикладів такої комуністичної діяльності в соціалістичному та капіталістичному суспільствах. Якщо розглядати ці та інші приклади з точки зору вираження волі щодо предметів праці, засобів праці та продукції, то йдеться про форми суспільної власності та суспільний характер присвоєння.

Наприклад, загальнодоступна база даних за своєю суттю вже не може перебувати не в суспільній власності, не кажучи вже про загальну систему управління економікою. Те саме стосується і вільного програмного забезпечення, і різноманітних сервісів, що передбачають спільне виробництво та використання. Існують також сервери для спільного використання ресурсів, які неефективно та невигідно використовувати на засадах приватної власності, наприклад, загальні бібліотеки та сховища для фільмів і музики, що передбачають, що учасники і отримують, і поширюють продукти на цих сервісах (але приклад із програмним забезпеченням і базами даних кращий, оскільки тут йдеться скоріше про виробництво, ніж про споживання).

З вищесказаного видно, що суспільна власність виникає внаслідок реального руху виробництва навіть там, де ще немає формального об’єднання засобів виробництва, в умовах капіталізму. Формальне об’єднання – захоплення політичної влади пролетаріатом і оволодіння ним засобами виробництва – дасть величезний простір для розвитку реальної суспільної власності. Зауважимо, що там, де є реальна, а не лише формальна суспільна власність, зникає підґрунтя для того, щоб намагатися привласнити собі частину суспільних ресурсів або, навпаки, розбазарювати їх, недбало ставитися до них, вважаючи їх «чиїмись».

У разі формальної суспільної власності необхідне зовнішнє регулювання, принаймні у формі моральної свідомості (фактично у вигляді закону та системи примусу, що забезпечує його виконання), щодо засобів, предметів і продуктів праці, яке унеможливлює подібні явища. У разі реальної суспільної власності потреба в такому регулюванні відпадає, оскільки культура ставлення до засобів виробництва стає просто умовою їх використання.
\end{document}
